% ============================================================================
% DEMO FRAGMENT -- new writing style for doc/manuscript.tex
% Merges, in prose, the current fragments:
%   - "The cell structure"  (subsec 136; subsubs 141,149,159: physical
%     properties / wavefront shaping / operational variables)
%   - "The light frame"     (419-422) + "Definitions" (447-468)
%   - "Interaction principles" (490-498) and the encounter description
%     "Superposing/Distinct bubbles" (502-568) condensed to one subsection;
%     the pair-formation rule tables move to an appendix.
% NOT standalone compilable; labels and macros are illustrative.  The goal is
% to show: one idea per paragraph, one notation table, claim->mechanism->role
% order, and explicit Postulate vs Derived wording.
% ============================================================================

\subsection{The cell state}
\label{subsec:cell-state-demo}

Every cell of the 4-dimensional lattice (three spatial torus directions plus
the internal dimension $W$) carries one fixed-size bit string.  That string
plays three distinct roles, and we state them once here so that the rest of
the model can be read without a glossary:

\begin{enumerate}
\item \textbf{Physical role.}  A six-bit charge word $ch$ is constant per
      $W$-layer and is partitioned into electric ($q$), weak ($w_1w_0$) and
      colour ($c_2c_1c_0$) sectors (Sect.~\ref{subsec:charges}).  Two derived
      bits, $pB=(pol_u>0)$ and $sB=(pol_v>0)$, mark the sign of the
      reconstructed polarisation pair and trigger the electric and magnetic
      interaction channels respectively.  The affinity $a$ labels the
      aggregate (island) a cell currently belongs to.
\item \textbf{Wavefront role.}  The squared radius $r^2$ and the integer
      radius $r$ are maintained incrementally from each moving source centre;
      a cell lies on the active spherical front iff $r$ equals the local
      frame radius $\mathrm{effective\_t}(t)$.  This radius grows by exactly
      one lattice cell per frame --- a postulate of the information clock,
      not an emergent speed (see Sect.~\ref{subsec:light-frame-demo}).
\item \textbf{Operational role.}  Two clocks coexist: the housekeeping tick
      $k$, incremented synchronously in every cell and used to schedule the
      FSM stages, and the local wavefront tick $t$, reset at each re-emitted
      source.  The vector $\mathbf{c}$ carries the pending translation, and
      $m$/$reloc$ the momentum and consumable displacement used by the
      inertial channel (Sect.~\ref{subsec:translation}).
\end{enumerate}

For reference, Table~\ref{tab:cell-state-demo} lists the same fields in one
place; it is a notation table, not a set of assumptions.

\begin{table}[ht]
\centering
\small
\begin{tabular}{lll}
\hline
Field & Role & Comment\\
\hline
$ch$      & physical & six-bit charge word, constant per $W$-layer\\
$pB,sB$   & physical & signs of reconstructed $(pol_u,pol_v)$; electric/magnetic channels\\
$a$       & physical & island/affinity label\\
$r^2,r$   & wavefront & incremental distance from source centre\\
$\mathit{active}$ & wavefront & true on the shell $r=\mathrm{effective\_t}(t)$\\
$t$, $f$  & wavefront & local clock and triangular breathing phase\\
$k$       & operational & global FSM stage counter\\
$\mathbf{c}$ & operational & pending translation offset\\
$m$, $reloc$ & operational & inertial momentum and consumable displacement\\
\hline
\end{tabular}
\caption{The cell bit string, grouped by role (single notation table).}
\label{tab:cell-state-demo}
\end{table}

Two clarifications are needed because the current text spreads them over
several short paragraphs.  First, all distances are toroidal and per-axis:
spatial edges are $L$ by default and the rectangular-tube allocator sets
them separately.  Second, the wavefront machinery is \emph{incremental}: it
never stores a full distance field, only the radius and its squared value
propagated through the six-neighbourhood with local square-boundary tests.

\subsection{Two clocks and the light frame}
\label{subsec:light-frame-demo}

The model separates \emph{timescales} from \emph{scheduling}.  An era is one
full expansion--contraction of a front; a light frame is one update cycle of
the whole lattice; an inertial re-emission occupies only a few housekeeping
ticks.  Every light frame advances the information clock by one unit: the
front moves exactly one lattice cell per frame by construction.  We therefore
do \emph{not} report a measured ``speed of one cell per tick'' as an emergent
result; that value is the definition of the clock, and the only empirical
question is whether the shell is complete at every radius.

Scheduling is handled by the global tick $k$, which walks through the five
stages of a frame: \textsc{encounter}, \textsc{diffusion}, \textsc{reloc},
\textsc{reissue} and \textsc{flood}.  Because $k$ is common to every cell,
encounter can compare the whole lattice along the internal dimension in one
sweep.  This comparison is \emph{non-local in scheduling}, and the text must
state plainly that no signalling claim follows from it: what needs proof is
that these scheduled reads never enter any cell's causal future in an
observable way, and the current implementation does not yet supply that proof
(its arbitration is host-side).  We therefore defer the ``no-signalling''
claim to a limitations paragraph rather than presenting it as established.

\subsection{Encounter and interaction principles}
\label{subsec:encounter-demo}

Encounter is the only place where two sources interact.  Its rules follow a
small number of principles, which the current text scatters across many short
subsections; we state them here and keep the detailed rule tables in the
appendix.

\begin{description}
\item[Reciprocity.]  Every interaction is evaluated twice, once with each
      source as the current cell; nothing writes the partner's draft.
\item[Identity before force.]  Equal-charge contacts first settle identity:
      an S/S pair elects a chief, a D/D pair promotes the lower address, and
      a K/K pair clashes so that the higher address becomes a delegate.  Only
      contacts that survive identity handling reach the charge/EM branches.
\item[Charge sectors.]  Complementary charges can annihilate into a
      pair (with a registered frequency unit); opposite-field singletons of
      the same sector merge under a dominant leader; same-field singletons
      repel one step.  The $s2B$ electroweak sieve gates only the
      electroweak channels and is gated by the active shell.
\item[Superposition vs distinctness.]  Equal pairs may superpose (blob);
      distinct-bubble rules separate same-sector, inter-sector and
      colour-neutral cases.  The colour-neutral rules exchange colour and
      reissue from the contact point; dispersion/singularization act at
      $t=RMAX/2$ between sectors.
\item[Light-matter channel.]  $pB$/$sB$ select electric/magnetic behaviour.
      When both bits are set the contact collapses ($kB=1$); when only one
      meets the other's surface the interaction is adiabatic: the sources
      exchange affinity and clock and relocate toward each other.
\end{description}

Diffusion then propagates the flags and contact information produced by
encounter synchronously through the six spatial neighbours, and translation
moves the source centres according to the accumulated $\mathbf{c}$ and the
inertial $reloc$ channel.  Because the encounter comparisons are scheduled
non-locally (above), the claim that the full cycle is a strictly local
finite-state rule is \emph{not} asserted for the present implementation; it
is listed as an open requirement in the limitations.

\subsection{Diffusion, translation and collapse in one reading}
\label{subsec:diffusion-demo}

Diffusion broadcasts the three things the next stage needs: the collapse bit
$kB$, the translation vector $\mathbf{c}$, and direction information.  A
collapse spreads to every bubble that shares the affinity of the collapsing
source, thereby marking the whole particle.  Independently, each cell copies
the first non-trivial $\mathbf{c}$ found among its six neighbours, so that by
the end of the stage every layer carries a single ready-to-use displacement;
$sB$-hunting increments $\mathbf{c}$ along any axis where a neighbour has
$hB=1$ and the local $sB=0$.  The phase $f$ is deliberately not diffused: it
is the local breathing value $\mathrm{effective\_t}(t)$ of each cell.

Translation consumes that displacement.  A bubble is copied from the
neighbour in the direction of $\mathbf{c}$ while the draft decrements
$\mathbf{c}$, so a source pattern shifts as a rigid block instead of being
rebuilt cell by cell; the inertial channel is separate and accumulates unit
steps into the consumable vector $reloc$ rather than a multi-cell jump.
Reissue then resets $t=0$ at the new source centre and forwards $f$ and the
active flag to the partner lattice; the old front continues until it fades,
and that leftover is what the text calls an \emph{orphan}.  Annihilation and
collapse complete the picture: a contact that sets $kB$ ends by re-emitting
the two bubbles from the contact point, which is how a front is reset rather
than teleported.  Exact slot numbers and the full transition function remain
in the algorithmic summary; this prose is meant as the map to that summary,
not as a replacement for it.

% --- what this demo merged, for the editor response --------------------------
% merged: subsubs "Physical/Wavefront/Operational" -> one subsection + one table
%         "Light frame" + "Definitions" paragraphs -> one subsection
%         "Interaction principles" + superposing/distinct fragments -> one
%           subsection with a principles list; pair tables -> appendix
%         "Diffusion"/"Translation"/"Annihilation and collapse" -> one subsection
% shortened: no-signalling claim downgraded; "measured 1 cell/tick" reframed
%            as clock definition (answers the circularity objection).

